% 本文件由 LaTeX 排版 Agent 根据有效 translation.json 自动生成。
% author=Ignatius of Antioch; work_id=0109; title=《依纳爵致斯米纳人书》

\chapter{《依纳爵致斯米纳人书》}

% structure: p0001 source=h1 target=omitted reason=page-title-already-rendered-by-parent

% structure: p0002 source=h2 target=section reason=relative-heading-rank; base=h2

\section{问候}

依纳爵，亦名背负天主者，致书于天主圣父及蒙爱者耶稣基督的教会——该教会赖仁慈获得各样恩赐，充满信德与爱德，一无所缺，堪受天主，并以圣德为饰：即位于亚细亚斯米纳的教会，愿藉无玷的圣神及天主的圣言，获享圆满的幸福。\par

% structure: p0004 source=h2 target=section reason=relative-heading-rank; base=h2

\section{第一章　为你们的信德感谢天主}

我光荣天主，即耶稣基督，因祂赐予你们如此智慧。因为我观察到，你们在坚定不移的信德上得以成全，彷佛被钉在我们主耶稣基督的十字架上，无论在肉身上还是精神上，并藉基督的血在爱德中立定，完全深信我们的主：按肉体说，祂确是出自达味的后裔\BibleRef{罗 1:3}，按天主的旨意和德能，祂是天主之子；祂确由童贞女所生，受若翰的洗礼，为要由祂满全一切义德\BibleRef{玛 3:15}；并在般雀比拉多和分封侯黑落德治下，确为我们被钉在十字架上，出于祂肉身的苦难。我们便是这果实的分享者，因祂神圣的苦难，祂为万世立起旗帜\BibleRef{依 5:26}，藉祂的复活，召集所有祂的圣洁和忠信者——无论是犹太人还是外邦人——进入祂教会同一身体内。\par

% structure: p0006 source=h2 target=section reason=relative-heading-rank; base=h2

\section{第二章　基督的真实苦难}

祂受了这一切，是为我们，使我们得救。祂确实受了苦，正如祂也确实复活了自己；不像某些不信者所主张的，祂只是似乎受苦，而他们自己也只是似乎为基督徒。他们怎样相信，事情也必怎样临到他们：当他们被剥去身体时，便成了纯粹的恶灵。\par

% structure: p0008 source=h2 target=section reason=relative-heading-rank; base=h2

\section{第三章　基督复活后仍具有身体}

因为我知道，祂复活后仍然具有肉身，并且我相信祂现在仍是如此。例如，当祂来到那些与伯多禄在一起的人那里时，对他们说：\Quote{你们摸一摸，看看我，就知道我不是一个无形的鬼魂。}他们立刻摸了祂，便相信了，既因祂的肉身，也因祂的精神而心服。为此，他们也轻看死亡，反而成为了死亡的征服者。祂复活后，又同他们一起吃喝，正如具有肉身者一样，虽然精神上祂与圣父合而为一。\par

% structure: p0010 source=h2 target=section reason=relative-heading-rank; base=h2

\section{第四章　谨防这些异端者}

亲爱的，我给你们这些指示，深知你们也持守同样的见解。但我预先防备你们，要提防那些披着人形的野兽；你们不仅不可接待他们，如有可能，甚至不要与他们相遇；但你们必须为他们祈祷，或许他们能得以悔改，尽管这非常困难。然而，我们的真生命耶稣基督有能力成就这事。但如果这些事只是表面上由我们的主所做，那么我也只是表面上被捆绑。我又为何将自己交付于死亡、烈火、刀剑、猛兽呢？其实，靠近刀剑者，就靠近天主；与野兽在一起者，就与天主为伴——只要他是因耶稣基督的名而如此。我经受这一切，为能同祂一同受苦\BibleRef{罗 8:17}，祂成了完美的人，在我内里坚固我\BibleRef{斐 4:13}。\par

% structure: p0012 source=h2 target=section reason=relative-heading-rank; base=h2

\section{第五章　他们危险的谬误}

有些人无知地否认祂，或者不如说被祂否认；他们是死亡的辩护者，而非真理的辩护者。这些人，既没有被先知说服，也没有被梅瑟法律说服，甚至直到今日也没有被福音说服，也没有被我们各自所受的苦难说服。因为他们对我们也有同样的看法。因为若有人称赞我，却亵渎我的主，不承认祂确实具有身体，这对我有何益处呢？但谁若不承认这一点，实际上就是完全否认了祂，被死亡所笼罩。然而，我不认为适宜写出这些人的名字，因为他们是不信者。是的，我绝不提及他们，直到他们悔改并回归于基督的苦难——这苦难就是我们的复活——的真正信仰。\par

% structure: p0014 source=h2 target=section reason=relative-heading-rank; base=h2

\section{第六章　不信基督之血者必受审判}

谁也不要自欺。无论是天上的事物、荣耀的天使、掌权者，可见的或不可见的，若不信基督的血，都将因此而招致定罪。\Quote{能领受的，就领受吧！}\BibleRef{玛 19:12}不要让高位使任何人自大；因为万有之上是信德和爱德，没有比它们更可贵的。但请注意那些对来到我们中间的基督恩宠持有不同意见的人，他们是何等反对天主的旨意。他们不重视爱德；不关心寡妇、孤儿、受压迫者，被捆绑的、自由的，饥饿的、口渴的。\par

% structure: p0016 source=h2 target=section reason=relative-heading-rank; base=h2

\section{第七章　让我们远离这样的异端者}

他们远离圣体圣事和祈祷，因为他们不承认圣体是我们救主耶稣基督的肉——这肉为我们的罪而受苦，且圣父因祂的仁慈使其复活。因此，那些反对天主这一恩赐的人，在争论中自取灭亡。但对他们来说，更好是尊重它，以便他们也能复活。所以，你们应当远离这样的人，不要在私下或公开谈论他们，而要留意先知，尤其是福音，其中基督的苦难已启示给我们，复活也已得到充分证明。但要避免一切分裂，因为它是邪恶的开端。\par

% structure: p0018 source=h2 target=section reason=relative-heading-rank; base=h2

\section{第八章　凡事不可没有主教}

你们都要跟随主教，如同耶稣基督跟随圣父；跟随长老团，如同跟随宗徒；尊敬执事，如同天主的制度。没有主教，谁也不可做任何与教会有关的事。只有由主教或他托付的人所举行的，才可视为合宜的圣体圣事。无论主教在哪里，民众也应聚集在那里；正如无论耶稣基督在哪里，那里就有天主教会。没有主教，无论是施洗或是举行爱筵，都是不合法的；但凡他所认可的，也必悦乐于天主，使所做的一切稳妥有效。\par

% structure: p0020 source=h2 target=section reason=relative-heading-rank; base=h2

\section{第九章　尊敬主教}

再者，我们理当恢复清醒的品行，趁着还有机会，向天主行悔改。尊敬天主和主教是好的。尊敬主教的人，已被天主所尊敬；凡未经主教知晓而行事的人，实际上是在事奉魔鬼。愿一切事物因恩宠而充盈于你们，因为你们是相称的。你们在一切事上安慰了我，耶稣基督也安慰了你们。你们无论在我缺席或在场时都爱了我。愿天主报答你们；为了祂的缘故，你们忍受一切，必将到达祂那里。\par

% structure: p0022 source=h2 target=section reason=relative-heading-rank; base=h2

\section{第十章　致谢他们的善行}

你们善待了菲罗和勒乌斯·阿伽托普斯，他们是我们天主基督的仆人，为了天主的缘故跟随了我，并因你们在各方面安慰了他们而向主感谢你们。这些事没有一件会失落于你们。愿我的灵和我的锁链为你们代求，你们没有轻视或羞于它们；我们的完美希望耶稣基督，也必不以你们为羞。\par

% structure: p0024 source=h2 target=section reason=relative-heading-rank; base=h2

\section{第十一章　请求他们派人去安提约基雅}

你们的祈祷已到达叙利亚安提约基雅的教会。我从那里带着锁链而来，最蒙天主悦纳，向众人致候；我不配从那里被称呼，因我是他们中最小的。然而，按天主的旨意，我被认为配得这份尊荣，并非我自觉有功劳，乃是因天主的恩宠，我愿这恩宠完全赐给我，好使因你们的祈祷，我能到达天主那里。为使你们的工作在地上和天上都得以成全，适合为天主的荣耀，你们的教会应选出一位可敬的使者；他前往叙利亚，可祝贺他们如今得享和平，恢复应有的伟大，并在他们中间重建应有的体制。在我看来，你们应当派一位代表带一封信函，与他们一同因按天主旨意所得的安宁而喜乐，并因你们的祈祷，他们现已抵达港湾。作为完美的人，你们也应追求完美的事。因为当你们渴望行善时，天主也乐意帮助你们。\par

% structure: p0026 source=h2 target=section reason=relative-heading-rank; base=h2

\section{第十二章　请安}

特洛阿的弟兄之爱向你们请安；我也从那里写信给你们，由布尔鲁斯带去，他与你们的弟兄厄弗所人同行，是你们派来与我一起的，他在一切事上安慰了我。愿人人都效法他，作为天主仆人的典范。恩宠必在一切事上报答他。我向你们极其可敬的主教、可敬的长老团、我的同仆执事们，并你们每一位，无论个人还是全体，因耶稣基督之名，并因祂的肉身和血、祂的苦难和复活——既是肉体的也是属灵的——以及与天主和你们的合一，向你们请安。愿恩宠、怜悯、平安和忍耐永远与你们同在！\par

% structure: p0028 source=h2 target=section reason=relative-heading-rank; base=h2

\section{第十三章　结语}

我向弟兄们的家属——他们的妻子儿女，以及被称为寡妇的贞女们——请安。我祈求你们因圣神的德能刚强。与我同在的菲罗向你们致候。我向塔维亚的家室请安，并祈求她在信德和爱德上，无论肉体还是属灵，都得以坚固。我向我亲爱的阿尔刻、无与伦比的达弗努斯、欧特克诺斯，以及各位按名请安。愿你们在天主的恩宠中安好。\par

\SourceRecord{Translated by Alexander Roberts and James Donaldson.；Ante-Nicene Fathers；Vol. 1.；Edited by Alexander Roberts, James Donaldson, and A. Cleveland Coxe.；Buffalo, NY: Christian Literature Publishing Co.,；1885.；<http://www.newadvent.org/fathers/0109.htm>.}
